% GENERATED FILE. Edit canonical lessons, then run npm run generate:books.
% canonical-source-hash: fnv1a64:2cf18cf7be5251c7
% canonical-lessons: TE-C237-polisu, TE-C237-draivaru, TE-C237-postumanu, TE-C237-kandaktaru, TE-C237-vakilu

\chapter{Police officer, Driver, Postman, Bus conductor, Lawyer}
\label{ch:te-police-officer-driver-postman-bus-conductor-lawyer}
\hlchaptermodality{237}

\begin{chapteropening}
\textbf{By the end of this chapter:} \emph{I can say a police officer, a driver, a postman, a bus conductor and a lawyer.}

Complete the last lesson of chapter 237: I can say a police officer, a driver, a postman, a bus conductor and a lawyer.
\end{chapteropening}

\section[\texorpdfstring{\te{పోలీసు} (pōlīsu)}{పోలీసు (pōlīsu)}]{\te{పోలీసు} (pōlīsu) — a police officer; the police}

\label{lesson:TE-C237-polisu}



\begin{quote}
Before the new one: say the Telugu for the government, then the Telugu for marks.
\end{quote}

\subsection*{You'll want to know: \te{పోలీసు}}
\textbf{\te{పోలీసు}} — \emph{pōlīsu} — \textquotedblleft{}a police officer; the police\textquotedblright{}.

\begin{grammarlens}[title={What you've built}]
One more word for this chapter.
\end{grammarlens}

\subsection*{Your turn}
\begin{itemize}
\raggedright
  \item \emph{Say it:} \emph{pōlīsu}
  \item \emph{Say it:} \emph{pōlīsu}, once more
  \item \emph{Say it:} say \emph{pōlīsu}
  \item \emph{Recall:} say the Telugu for the government, then the Telugu for marks, then say \emph{pōlīsu} again
\end{itemize}

\begin{tcolorbox}[breakable,colback=teal!4,colframe=teal!35!black,title={Before you move on}]
What does \te{పోలీసు} mean? (\textquotedblleft{}A police officer; the police\textquotedblright{}.) Say it once more.
\end{tcolorbox}
\section[\texorpdfstring{\te{డ్రైవరు} (ḍraivaru)}{డ్రైవరు (ḍraivaru)}]{\te{డ్రైవరు} (ḍraivaru) — a driver}

\label{lesson:TE-C237-draivaru}



\begin{quote}
Before the new one: say the Telugu for marks, then the Telugu for a police officer.
\end{quote}

\subsection*{You'll want to know: \te{డ్రైవరు}}
\textbf{\te{డ్రైవరు}} — \emph{ḍraivaru} — \textquotedblleft{}a driver\textquotedblright{}.

\begin{grammarlens}[title={What you've built}]
One more word for this chapter.
\end{grammarlens}

\subsection*{Your turn}
\begin{itemize}
\raggedright
  \item \emph{Say it:} \emph{ḍraivaru}
  \item \emph{Say it:} \emph{ḍraivaru}, once more
  \item \emph{Say it:} say \emph{ḍraivaru}
  \item \emph{Recall:} say the Telugu for marks, then the Telugu for a police officer, then say \emph{ḍraivaru} again
\end{itemize}

\begin{tcolorbox}[breakable,colback=teal!4,colframe=teal!35!black,title={Before you move on}]
What does \te{డ్రైవరు} mean? (\textquotedblleft{}A driver\textquotedblright{}.) Say it once more.
\end{tcolorbox}
\section[\texorpdfstring{\te{పోస్టుమాను} (pōsṭumānu)}{పోస్టుమాను (pōsṭumānu)}]{\te{పోస్టుమాను} (pōsṭumānu) — a postman}

\label{lesson:TE-C237-postumanu}



\begin{quote}
Before the new one: say the Telugu for a police officer, then the Telugu for a driver.
\end{quote}

\subsection*{You'll want to know: \te{పోస్టుమాను}}
\textbf{\te{పోస్టుమాను}} — \emph{pōsṭumānu} — \textquotedblleft{}a postman\textquotedblright{}.

\begin{grammarlens}[title={What you've built}]
One more word for this chapter.
\end{grammarlens}

\subsection*{Your turn}
\begin{itemize}
\raggedright
  \item \emph{Say it:} \emph{pōsṭumānu}
  \item \emph{Say it:} \emph{pōsṭumānu}, once more
  \item \emph{Say it:} say \emph{pōsṭumānu}
  \item \emph{Recall:} say the Telugu for a police officer, then the Telugu for a driver, then say \emph{pōsṭumānu} again
\end{itemize}

\begin{tcolorbox}[breakable,colback=teal!4,colframe=teal!35!black,title={Before you move on}]
What does \te{పోస్టుమాను} mean? (\textquotedblleft{}A postman\textquotedblright{}.) Say it once more.
\end{tcolorbox}
\section[\texorpdfstring{\te{కండక్టరు} (kaṇḍakṭaru)}{కండక్టరు (kaṇḍakṭaru)}]{\te{కండక్టరు} (kaṇḍakṭaru) — a bus conductor}

\label{lesson:TE-C237-kandaktaru}



\begin{quote}
Before the new one: say the Telugu for a driver, then the Telugu for a postman.
\end{quote}

\subsection*{You'll want to know: \te{కండక్టరు}}
\textbf{\te{కండక్టరు}} — \emph{kaṇḍakṭaru} — \textquotedblleft{}a bus conductor\textquotedblright{}.

\begin{grammarlens}[title={What you've built}]
One more word for this chapter.
\end{grammarlens}

\subsection*{Your turn}
\begin{itemize}
\raggedright
  \item \emph{Say it:} \emph{kaṇḍakṭaru}
  \item \emph{Say it:} \emph{kaṇḍakṭaru}, once more
  \item \emph{Say it:} say \emph{kaṇḍakṭaru}
  \item \emph{Recall:} say the Telugu for a driver, then the Telugu for a postman, then say \emph{kaṇḍakṭaru} again
\end{itemize}

\begin{tcolorbox}[breakable,colback=teal!4,colframe=teal!35!black,title={Before you move on}]
What does \te{కండక్టరు} mean? (\textquotedblleft{}A bus conductor\textquotedblright{}.) Say it once more.
\end{tcolorbox}
\section[\texorpdfstring{\te{వకీలు} (vakīlu)}{వకీలు (vakīlu)}]{\te{వకీలు} (vakīlu) — a lawyer}

\label{lesson:TE-C237-vakilu}



\begin{quote}
Before the new one: say the Telugu for a postman, then the Telugu for a bus conductor.
\end{quote}

\subsection*{You'll want to know: \te{వకీలు}}
\textbf{\te{వకీలు}} — \emph{vakīlu} — \textquotedblleft{}a lawyer\textquotedblright{}.

\begin{grammarlens}[title={What you've built}]
One more word for this chapter.
\end{grammarlens}

\subsection*{Your turn}
\begin{itemize}
\raggedright
  \item \emph{Say it:} \emph{vakīlu}
  \item \emph{Say it:} \emph{vakīlu}, once more
  \item \emph{Say it:} say \emph{vakīlu}
  \item \emph{Recall:} say the Telugu for a postman, then the Telugu for a bus conductor, then say \emph{vakīlu} again
\end{itemize}

\begin{tcolorbox}[breakable,colback=teal!4,colframe=teal!35!black,title={Before you move on}]
What does \te{వకీలు} mean? (\textquotedblleft{}A lawyer\textquotedblright{}.) Say it once more.
\end{tcolorbox}
